% =============================================================================
% Real-city experiment design (Pittsburgh / Chicago / Manhattan / Los Angeles)
%
% Drop into your Overleaf project and % =============================================================================
% Real-city experiment design (Pittsburgh / Chicago / Manhattan / Los Angeles)
%
% Drop into your Overleaf project and % =============================================================================
% Real-city experiment design (Pittsburgh / Chicago / Manhattan / Los Angeles)
%
% Drop into your Overleaf project and % =============================================================================
% Real-city experiment design (Pittsburgh / Chicago / Manhattan / Los Angeles)
%
% Drop into your Overleaf project and \input{real_city_experiment} where
% needed. Self-contained: no other doc cross-references required.
% Requires: amsmath, booktabs, array, hyperref (optional).
% =============================================================================

\subsection{Real-City Experiment Design}\label{sec:real-city-design}

This section documents the design of the real-city baseline sweep
described in \S\ref{sec:multi-city}. The goal is to compare the seven
service-design baselines on the same regime grid as the synthetic
$5\times5$ run, but on \emph{actual} U.S.\ census block groups (BGs)
linked together by \emph{actual} street networks --- preserving
real demand density (population) and real driving distances.

% -----------------------------------------------------------------------
\subsubsection{Cities and their road topologies}

Four cities were chosen to span the major U.S.\ road-network archetypes:

\begin{center}
\small
\begin{tabular}{lll}
\toprule
City & Focal area / topology & BG suburb \\
\midrule
Pittsburgh, PA & Downtown + Oakland; hilly, river-cut, irregular grid &
   Allegheny Co.\ (FIPS 42003) \\
Chicago, IL    & Loop + River North; strict cardinal grid &
   Cook Co.\ (FIPS 17031) \\
Manhattan, NY  & Midtown; avenue-and-street grid plus Broadway &
   New York Co.\ (FIPS 36061) \\
Los Angeles, CA & DTLA; sprawled grid bisected by freeways &
   LA Co.\ (FIPS 06037) \\
\bottomrule
\end{tabular}
\end{center}

For each city we keep the $25$ block groups whose centroids are
closest to the centroid of a focal $\sim 5\,$km bbox. Twenty-five BGs
match the synthetic $5\times5$ instance and keep every per-baseline
sweep tractable on a workstation; the centermost selection avoids
boundary BGs whose road network is partially clipped.

% -----------------------------------------------------------------------
\subsubsection{Data sources}

\begin{description}
  \item[Block-group polygons] U.S.\ Census Bureau TIGER 2023
    shapefiles (one file per state, downloaded via
    \texttt{www2.census.gov/geo/tiger/TIGER2023/BG}). Reprojected to
    EPSG:2163 for area + Euclidean distance computations.
  \item[Population per BG] American Community Survey 5-year 2023
    table B01003 (total population), fetched per BG via the public
    Census Data API. Saved as a per-city JSON
    \verb|data/city_bg_populations/<city>.json|.
  \item[Road network] OpenStreetMap drive-network graph fetched via
    OSMnx 2.x for each city's bbox (with a $0.005^\circ$ pad).
    Cached on disk under \verb|cache/road_graph_*.pkl|.
\end{description}

The end-to-end fetch takes a few minutes per city; helper scripts
(\verb|build_city_geoid_lists.py|, \verb|fetch_city_populations.py|)
handle this once and write the artefacts to disk for reproducible
sweeps.

% -----------------------------------------------------------------------
\subsubsection{Distances: road network + MDS embedding}

The repository historically used Euclidean centroid-to-centroid
distances throughout. For real cities we replace these with
\textbf{road-network shortest-path distances} on three layers:

\paragraph{1. Centroid pair distances (\texttt{geodata.get\_dist}).}
After OSMnx fetches the drive graph, every BG centroid is mapped to
its nearest road-network node and an all-pairs shortest-path matrix
$D_{\mathrm{road}} \in \mathbb{R}^{N \times N}$ is precomputed
(\verb|GeoData.attach_road_distances|). Once attached,
\verb|GeoData.get_dist(b_1, b_2)| returns the road km from
$D_{\mathrm{road}}$, so any baseline that already calls
\verb|get_dist| (e.g.\ Joint-Nom/DRO and SP-Lit's linehaul
$K_i$) automatically becomes road-aware without code changes.

\paragraph{2. MDS-embedded positions for routing helpers.}
Several baselines (Multi-FR, TP-Lit, VCC) compute pairwise
distances via \verb|np.linalg.norm| over centroid coordinates
in EPSG:2163 metres. Replacing every call site with matrix
look-ups would be invasive. Instead, when the road-distance matrix
is attached we compute a \textbf{2D MDS embedding}
$\Phi : \mathrm{BG} \to \mathbb{R}^2$ such that
$\| \Phi(b_1) - \Phi(b_2) \|_2 \approx D_{\mathrm{road}}(b_1, b_2)$
and expose it as \verb|geodata.pos_road_km|. The embedding is
computed with scikit-learn's \verb|sklearn.manifold.MDS|
(precomputed dissimilarity, $n=2$, $4$ initialisations,
$300$ iterations).

A new helper \verb|_get_pos_road_aware(geodata, b)| returns
\verb|pos_road_km[b]| when available and falls back to the
Euclidean \verb|pos[b]| otherwise. We swapped this in for the
position lookups inside Multi-FR's line library, TP-Lit's ADP+VRP
candidate locations, and VCC's depot/dest/start positions. The
existing $L^2$-norm logic then operates on road-distance-consistent
geometry.

The MDS embedding is approximate --- general road-distance metrics
are not exactly Euclidean-embeddable in two dimensions --- but the
stress is moderate (typical pairwise error $\le 30\%$ on these 25-BG
instances) and the embedding propagates road awareness through
hundreds of $L^2$-norm call sites without bespoke refactors.

\paragraph{3. Continuous-sample distances in OD.}
OD samples customer destinations as continuous points within block
polygons, then routes \verb|depot -> destination -> depot|. We
\textbf{centroid-snap} each sampled point to its containing BG and
look up the road km from $D_{\mathrm{road}}$. This is faster than
on-the-fly road queries per sample and accurate to within one
centroid-radius ($\approx 100$--$300$\,m for our 25-BG areas).

\paragraph{Summary of distance use per baseline.}
\begin{center}
\small
\begin{tabular}{l p{4.6cm} p{4.6cm}}
\toprule
Baseline & Optimisation distances & Evaluation distances \\
\midrule
Joint-Nom / Joint-DRO &
   BHH $\alpha_i = \beta\sqrt{\Lambda}\sum_b \sqrt{p_b A_b}$;
   $K_i$ via road km ($\verb|get_dist|$ override). &
   In-district NN-TSP on road km via
   $\verb|road_network.get_dist_matrix_km|$. \\
SP-Lit (Carlsson) &
   Same as Joint; ham-sandwich bisection on Euclidean grid in
   EPSG:2163 metres (BG centroids unchanged). &
   Same as Joint. \\
Multi-FR / Multi-FR-Detour &
   Angular partition + 2-opt + subpath LP on road-aware MDS
   positions $\Phi(b)$ (Euclidean over $\Phi$). &
   Custom Benders simulate; Euclidean over $\Phi$. \\
TP-Lit (Liu) &
   ADP+VRP cost matrices via Euclidean over $\Phi$
   ($n_{\mathrm{candidates}} = 0$ so candidates are real BG
   centroids). &
   Custom simulate; Euclidean over $\Phi$. \\
VCC &
   Dial-a-ride routing via Euclidean over $\Phi$. &
   Same. Continuous walk distances within blocks remain Euclidean. \\
OD &
   Single trip $\verb|depot -> dest -> depot|$. &
   Centroid-snap road km from $D_{\mathrm{road}}$. \\
\bottomrule
\end{tabular}
\end{center}

% -----------------------------------------------------------------------
\subsubsection{Demand}

Demand at each BG is proportional to its ACS 2023 total population
(table B01003).
Per-block probabilities are normalised to sum to one over the 25-BG
subset:

\begin{equation*}
  p_b = \frac{\mathrm{pop}_b}{\sum_{b' \in \mathrm{subset}} \mathrm{pop}_{b'}}.
\end{equation*}

Total subset populations vary across cities (Pittsburgh: 24{,}604;
Chicago: 63{,}522; Manhattan: 26{,}326; Los Angeles: 37{,}428),
which alone introduces meaningful per-block density differences ---
Manhattan's 25 BGs are roughly half the area of Pittsburgh's but
have similar total population. The system arrival rate $\Lambda$
is swept independently of population total; the population
distribution only sets the spatial pattern $p_b$.

% -----------------------------------------------------------------------
\subsubsection{Sweep grid (regime view)}

Identical to the synthetic $5\times5$ design (\S\ref{sec:regime-view})
so cross-instance comparison is direct:

\begin{itemize}
  \item $\Lambda \in \{50, 150, 300\}$ pax/h
  \item $w_r \in \{1, 3, 10\}$
  \item $w_v = 1.0$ (pinned)
  \item $c_f = 4.0$ \$/veh-h (5-yr fleet amortisation)
  \item $K = 3$ baseline (used by Multi-FR's fleet fall-back; partition methods sweep $K \in \{2,...,6\}$ internally)
\end{itemize}

For each $(\Lambda, w_r)$ point, every baseline runs its
internal hyperparameter sweep
(\S\ref{sec:sweep-params}: K for partition, $T_{\mathrm{fixed}}$ for
SP-Lit, fleet for OD/VCC, etc.) and reports the operating point
with the lowest aggregate cost
$C_{\mathrm{provider}} + w_r \cdot \bar{u} \cdot \Lambda$.

% -----------------------------------------------------------------------
\subsubsection{Reproducibility}

Three scripts implement the pipeline end to end:

\begin{enumerate}
  \item \texttt{scripts/experiments/build\_city\_geoid\_lists.py} \\
    Filter each state's TIGER BG shapefile by the city's focal bbox,
    keep the $25$ centermost BGs, save GEOID lists to
    \verb|data/city_bg_lists/<city>.txt|.
  \item \texttt{scripts/experiments/fetch\_city\_populations.py} \\
    Query the Census API for ACS 2023 B01003 per BG, save
    per-city JSON to \verb|data/city_bg_populations/<city>.json|.
  \item \texttt{scripts/experiments/run\_per\_city\_full.sh} \\
    Run the regime sweep across all 4 cities with road-aware
    optimisation, population-weighted demand, and all 7 baselines
    (Multi-FR + TP-Lit included). Resumable via per-setting JSON
    cache.
\end{enumerate}

The plotting helpers
\verb|scripts/plots/multi_city_designs_map.py| and the city-aware
extension of \verb|rebuild_latest_sweeps.py| produce the
4-panel scatter cloud and the $4 \times N_{\mathrm{baseline}}$
design map shown in Fig.~\ref{fig:cities-cloud} and
Fig.~\ref{fig:cities-designs}.

% -----------------------------------------------------------------------
\subsubsection{Known limitations}

\begin{itemize}
  \item \textbf{MDS approximation error.} Road-distance metrics are
    not exactly Euclidean-embeddable in 2D; pairwise distances
    after embedding can deviate by up to $\sim 30\%$ on irregular
    networks (Pittsburgh). A stricter alternative is to refactor
    each baseline's routing helpers to consume an explicit
    distance-matrix function; this is left for future work.
  \item \textbf{Continuous-point road distance in VCC.} VCC's
    inner dial-a-ride routing operates on continuous coordinates
    (interpolated meeting points). Those interior moves remain
    Euclidean over MDS positions; only depot, destination, and
    block-anchor distances are road-aware.
  \item \textbf{Demand temporality.} ACS B01003 is total residential
    population, not time-varying ridership. A more realistic
    deployment study would weight by trip-generation tables
    (e.g.\ ACS B08303 commuting) at peak hours. Population-weighted
    demand is a stand-in for cross-city methodology validation.
  \item \textbf{Single bbox per city.} Each city is represented by
    a single $\sim 5\,$km focal area. Multi-bbox runs (e.g.\
    Pittsburgh downtown + suburbs) would test whether topology
    rankings hold across intra-city heterogeneity.
\end{itemize}
 where
% needed. Self-contained: no other doc cross-references required.
% Requires: amsmath, booktabs, array, hyperref (optional).
% =============================================================================

\subsection{Real-City Experiment Design}\label{sec:real-city-design}

This section documents the design of the real-city baseline sweep
described in \S\ref{sec:multi-city}. The goal is to compare the seven
service-design baselines on the same regime grid as the synthetic
$5\times5$ run, but on \emph{actual} U.S.\ census block groups (BGs)
linked together by \emph{actual} street networks --- preserving
real demand density (population) and real driving distances.

% -----------------------------------------------------------------------
\subsubsection{Cities and their road topologies}

Four cities were chosen to span the major U.S.\ road-network archetypes:

\begin{center}
\small
\begin{tabular}{lll}
\toprule
City & Focal area / topology & BG suburb \\
\midrule
Pittsburgh, PA & Downtown + Oakland; hilly, river-cut, irregular grid &
   Allegheny Co.\ (FIPS 42003) \\
Chicago, IL    & Loop + River North; strict cardinal grid &
   Cook Co.\ (FIPS 17031) \\
Manhattan, NY  & Midtown; avenue-and-street grid plus Broadway &
   New York Co.\ (FIPS 36061) \\
Los Angeles, CA & DTLA; sprawled grid bisected by freeways &
   LA Co.\ (FIPS 06037) \\
\bottomrule
\end{tabular}
\end{center}

For each city we keep the $25$ block groups whose centroids are
closest to the centroid of a focal $\sim 5\,$km bbox. Twenty-five BGs
match the synthetic $5\times5$ instance and keep every per-baseline
sweep tractable on a workstation; the centermost selection avoids
boundary BGs whose road network is partially clipped.

% -----------------------------------------------------------------------
\subsubsection{Data sources}

\begin{description}
  \item[Block-group polygons] U.S.\ Census Bureau TIGER 2023
    shapefiles (one file per state, downloaded via
    \texttt{www2.census.gov/geo/tiger/TIGER2023/BG}). Reprojected to
    EPSG:2163 for area + Euclidean distance computations.
  \item[Population per BG] American Community Survey 5-year 2023
    table B01003 (total population), fetched per BG via the public
    Census Data API. Saved as a per-city JSON
    \verb|data/city_bg_populations/<city>.json|.
  \item[Road network] OpenStreetMap drive-network graph fetched via
    OSMnx 2.x for each city's bbox (with a $0.005^\circ$ pad).
    Cached on disk under \verb|cache/road_graph_*.pkl|.
\end{description}

The end-to-end fetch takes a few minutes per city; helper scripts
(\verb|build_city_geoid_lists.py|, \verb|fetch_city_populations.py|)
handle this once and write the artefacts to disk for reproducible
sweeps.

% -----------------------------------------------------------------------
\subsubsection{Distances: road network + MDS embedding}

The repository historically used Euclidean centroid-to-centroid
distances throughout. For real cities we replace these with
\textbf{road-network shortest-path distances} on three layers:

\paragraph{1. Centroid pair distances (\texttt{geodata.get\_dist}).}
After OSMnx fetches the drive graph, every BG centroid is mapped to
its nearest road-network node and an all-pairs shortest-path matrix
$D_{\mathrm{road}} \in \mathbb{R}^{N \times N}$ is precomputed
(\verb|GeoData.attach_road_distances|). Once attached,
\verb|GeoData.get_dist(b_1, b_2)| returns the road km from
$D_{\mathrm{road}}$, so any baseline that already calls
\verb|get_dist| (e.g.\ Joint-Nom/DRO and SP-Lit's linehaul
$K_i$) automatically becomes road-aware without code changes.

\paragraph{2. MDS-embedded positions for routing helpers.}
Several baselines (Multi-FR, TP-Lit, VCC) compute pairwise
distances via \verb|np.linalg.norm| over centroid coordinates
in EPSG:2163 metres. Replacing every call site with matrix
look-ups would be invasive. Instead, when the road-distance matrix
is attached we compute a \textbf{2D MDS embedding}
$\Phi : \mathrm{BG} \to \mathbb{R}^2$ such that
$\| \Phi(b_1) - \Phi(b_2) \|_2 \approx D_{\mathrm{road}}(b_1, b_2)$
and expose it as \verb|geodata.pos_road_km|. The embedding is
computed with scikit-learn's \verb|sklearn.manifold.MDS|
(precomputed dissimilarity, $n=2$, $4$ initialisations,
$300$ iterations).

A new helper \verb|_get_pos_road_aware(geodata, b)| returns
\verb|pos_road_km[b]| when available and falls back to the
Euclidean \verb|pos[b]| otherwise. We swapped this in for the
position lookups inside Multi-FR's line library, TP-Lit's ADP+VRP
candidate locations, and VCC's depot/dest/start positions. The
existing $L^2$-norm logic then operates on road-distance-consistent
geometry.

The MDS embedding is approximate --- general road-distance metrics
are not exactly Euclidean-embeddable in two dimensions --- but the
stress is moderate (typical pairwise error $\le 30\%$ on these 25-BG
instances) and the embedding propagates road awareness through
hundreds of $L^2$-norm call sites without bespoke refactors.

\paragraph{3. Continuous-sample distances in OD.}
OD samples customer destinations as continuous points within block
polygons, then routes \verb|depot -> destination -> depot|. We
\textbf{centroid-snap} each sampled point to its containing BG and
look up the road km from $D_{\mathrm{road}}$. This is faster than
on-the-fly road queries per sample and accurate to within one
centroid-radius ($\approx 100$--$300$\,m for our 25-BG areas).

\paragraph{Summary of distance use per baseline.}
\begin{center}
\small
\begin{tabular}{l p{4.6cm} p{4.6cm}}
\toprule
Baseline & Optimisation distances & Evaluation distances \\
\midrule
Joint-Nom / Joint-DRO &
   BHH $\alpha_i = \beta\sqrt{\Lambda}\sum_b \sqrt{p_b A_b}$;
   $K_i$ via road km ($\verb|get_dist|$ override). &
   In-district NN-TSP on road km via
   $\verb|road_network.get_dist_matrix_km|$. \\
SP-Lit (Carlsson) &
   Same as Joint; ham-sandwich bisection on Euclidean grid in
   EPSG:2163 metres (BG centroids unchanged). &
   Same as Joint. \\
Multi-FR / Multi-FR-Detour &
   Angular partition + 2-opt + subpath LP on road-aware MDS
   positions $\Phi(b)$ (Euclidean over $\Phi$). &
   Custom Benders simulate; Euclidean over $\Phi$. \\
TP-Lit (Liu) &
   ADP+VRP cost matrices via Euclidean over $\Phi$
   ($n_{\mathrm{candidates}} = 0$ so candidates are real BG
   centroids). &
   Custom simulate; Euclidean over $\Phi$. \\
VCC &
   Dial-a-ride routing via Euclidean over $\Phi$. &
   Same. Continuous walk distances within blocks remain Euclidean. \\
OD &
   Single trip $\verb|depot -> dest -> depot|$. &
   Centroid-snap road km from $D_{\mathrm{road}}$. \\
\bottomrule
\end{tabular}
\end{center}

% -----------------------------------------------------------------------
\subsubsection{Demand}

Demand at each BG is proportional to its ACS 2023 total population
(table B01003).
Per-block probabilities are normalised to sum to one over the 25-BG
subset:

\begin{equation*}
  p_b = \frac{\mathrm{pop}_b}{\sum_{b' \in \mathrm{subset}} \mathrm{pop}_{b'}}.
\end{equation*}

Total subset populations vary across cities (Pittsburgh: 24{,}604;
Chicago: 63{,}522; Manhattan: 26{,}326; Los Angeles: 37{,}428),
which alone introduces meaningful per-block density differences ---
Manhattan's 25 BGs are roughly half the area of Pittsburgh's but
have similar total population. The system arrival rate $\Lambda$
is swept independently of population total; the population
distribution only sets the spatial pattern $p_b$.

% -----------------------------------------------------------------------
\subsubsection{Sweep grid (regime view)}

Identical to the synthetic $5\times5$ design (\S\ref{sec:regime-view})
so cross-instance comparison is direct:

\begin{itemize}
  \item $\Lambda \in \{50, 150, 300\}$ pax/h
  \item $w_r \in \{1, 3, 10\}$
  \item $w_v = 1.0$ (pinned)
  \item $c_f = 4.0$ \$/veh-h (5-yr fleet amortisation)
  \item $K = 3$ baseline (used by Multi-FR's fleet fall-back; partition methods sweep $K \in \{2,...,6\}$ internally)
\end{itemize}

For each $(\Lambda, w_r)$ point, every baseline runs its
internal hyperparameter sweep
(\S\ref{sec:sweep-params}: K for partition, $T_{\mathrm{fixed}}$ for
SP-Lit, fleet for OD/VCC, etc.) and reports the operating point
with the lowest aggregate cost
$C_{\mathrm{provider}} + w_r \cdot \bar{u} \cdot \Lambda$.

% -----------------------------------------------------------------------
\subsubsection{Reproducibility}

Three scripts implement the pipeline end to end:

\begin{enumerate}
  \item \texttt{scripts/experiments/build\_city\_geoid\_lists.py} \\
    Filter each state's TIGER BG shapefile by the city's focal bbox,
    keep the $25$ centermost BGs, save GEOID lists to
    \verb|data/city_bg_lists/<city>.txt|.
  \item \texttt{scripts/experiments/fetch\_city\_populations.py} \\
    Query the Census API for ACS 2023 B01003 per BG, save
    per-city JSON to \verb|data/city_bg_populations/<city>.json|.
  \item \texttt{scripts/experiments/run\_per\_city\_full.sh} \\
    Run the regime sweep across all 4 cities with road-aware
    optimisation, population-weighted demand, and all 7 baselines
    (Multi-FR + TP-Lit included). Resumable via per-setting JSON
    cache.
\end{enumerate}

The plotting helpers
\verb|scripts/plots/multi_city_designs_map.py| and the city-aware
extension of \verb|rebuild_latest_sweeps.py| produce the
4-panel scatter cloud and the $4 \times N_{\mathrm{baseline}}$
design map shown in Fig.~\ref{fig:cities-cloud} and
Fig.~\ref{fig:cities-designs}.

% -----------------------------------------------------------------------
\subsubsection{Known limitations}

\begin{itemize}
  \item \textbf{MDS approximation error.} Road-distance metrics are
    not exactly Euclidean-embeddable in 2D; pairwise distances
    after embedding can deviate by up to $\sim 30\%$ on irregular
    networks (Pittsburgh). A stricter alternative is to refactor
    each baseline's routing helpers to consume an explicit
    distance-matrix function; this is left for future work.
  \item \textbf{Continuous-point road distance in VCC.} VCC's
    inner dial-a-ride routing operates on continuous coordinates
    (interpolated meeting points). Those interior moves remain
    Euclidean over MDS positions; only depot, destination, and
    block-anchor distances are road-aware.
  \item \textbf{Demand temporality.} ACS B01003 is total residential
    population, not time-varying ridership. A more realistic
    deployment study would weight by trip-generation tables
    (e.g.\ ACS B08303 commuting) at peak hours. Population-weighted
    demand is a stand-in for cross-city methodology validation.
  \item \textbf{Single bbox per city.} Each city is represented by
    a single $\sim 5\,$km focal area. Multi-bbox runs (e.g.\
    Pittsburgh downtown + suburbs) would test whether topology
    rankings hold across intra-city heterogeneity.
\end{itemize}
 where
% needed. Self-contained: no other doc cross-references required.
% Requires: amsmath, booktabs, array, hyperref (optional).
% =============================================================================

\subsection{Real-City Experiment Design}\label{sec:real-city-design}

This section documents the design of the real-city baseline sweep
described in \S\ref{sec:multi-city}. The goal is to compare the seven
service-design baselines on the same regime grid as the synthetic
$5\times5$ run, but on \emph{actual} U.S.\ census block groups (BGs)
linked together by \emph{actual} street networks --- preserving
real demand density (population) and real driving distances.

% -----------------------------------------------------------------------
\subsubsection{Cities and their road topologies}

Four cities were chosen to span the major U.S.\ road-network archetypes:

\begin{center}
\small
\begin{tabular}{lll}
\toprule
City & Focal area / topology & BG suburb \\
\midrule
Pittsburgh, PA & Downtown + Oakland; hilly, river-cut, irregular grid &
   Allegheny Co.\ (FIPS 42003) \\
Chicago, IL    & Loop + River North; strict cardinal grid &
   Cook Co.\ (FIPS 17031) \\
Manhattan, NY  & Midtown; avenue-and-street grid plus Broadway &
   New York Co.\ (FIPS 36061) \\
Los Angeles, CA & DTLA; sprawled grid bisected by freeways &
   LA Co.\ (FIPS 06037) \\
\bottomrule
\end{tabular}
\end{center}

For each city we keep the $25$ block groups whose centroids are
closest to the centroid of a focal $\sim 5\,$km bbox. Twenty-five BGs
match the synthetic $5\times5$ instance and keep every per-baseline
sweep tractable on a workstation; the centermost selection avoids
boundary BGs whose road network is partially clipped.

% -----------------------------------------------------------------------
\subsubsection{Data sources}

\begin{description}
  \item[Block-group polygons] U.S.\ Census Bureau TIGER 2023
    shapefiles (one file per state, downloaded via
    \texttt{www2.census.gov/geo/tiger/TIGER2023/BG}). Reprojected to
    EPSG:2163 for area + Euclidean distance computations.
  \item[Population per BG] American Community Survey 5-year 2023
    table B01003 (total population), fetched per BG via the public
    Census Data API. Saved as a per-city JSON
    \verb|data/city_bg_populations/<city>.json|.
  \item[Road network] OpenStreetMap drive-network graph fetched via
    OSMnx 2.x for each city's bbox (with a $0.005^\circ$ pad).
    Cached on disk under \verb|cache/road_graph_*.pkl|.
\end{description}

The end-to-end fetch takes a few minutes per city; helper scripts
(\verb|build_city_geoid_lists.py|, \verb|fetch_city_populations.py|)
handle this once and write the artefacts to disk for reproducible
sweeps.

% -----------------------------------------------------------------------
\subsubsection{Distances: road network + MDS embedding}

The repository historically used Euclidean centroid-to-centroid
distances throughout. For real cities we replace these with
\textbf{road-network shortest-path distances} on three layers:

\paragraph{1. Centroid pair distances (\texttt{geodata.get\_dist}).}
After OSMnx fetches the drive graph, every BG centroid is mapped to
its nearest road-network node and an all-pairs shortest-path matrix
$D_{\mathrm{road}} \in \mathbb{R}^{N \times N}$ is precomputed
(\verb|GeoData.attach_road_distances|). Once attached,
\verb|GeoData.get_dist(b_1, b_2)| returns the road km from
$D_{\mathrm{road}}$, so any baseline that already calls
\verb|get_dist| (e.g.\ Joint-Nom/DRO and SP-Lit's linehaul
$K_i$) automatically becomes road-aware without code changes.

\paragraph{2. MDS-embedded positions for routing helpers.}
Several baselines (Multi-FR, TP-Lit, VCC) compute pairwise
distances via \verb|np.linalg.norm| over centroid coordinates
in EPSG:2163 metres. Replacing every call site with matrix
look-ups would be invasive. Instead, when the road-distance matrix
is attached we compute a \textbf{2D MDS embedding}
$\Phi : \mathrm{BG} \to \mathbb{R}^2$ such that
$\| \Phi(b_1) - \Phi(b_2) \|_2 \approx D_{\mathrm{road}}(b_1, b_2)$
and expose it as \verb|geodata.pos_road_km|. The embedding is
computed with scikit-learn's \verb|sklearn.manifold.MDS|
(precomputed dissimilarity, $n=2$, $4$ initialisations,
$300$ iterations).

A new helper \verb|_get_pos_road_aware(geodata, b)| returns
\verb|pos_road_km[b]| when available and falls back to the
Euclidean \verb|pos[b]| otherwise. We swapped this in for the
position lookups inside Multi-FR's line library, TP-Lit's ADP+VRP
candidate locations, and VCC's depot/dest/start positions. The
existing $L^2$-norm logic then operates on road-distance-consistent
geometry.

The MDS embedding is approximate --- general road-distance metrics
are not exactly Euclidean-embeddable in two dimensions --- but the
stress is moderate (typical pairwise error $\le 30\%$ on these 25-BG
instances) and the embedding propagates road awareness through
hundreds of $L^2$-norm call sites without bespoke refactors.

\paragraph{3. Continuous-sample distances in OD.}
OD samples customer destinations as continuous points within block
polygons, then routes \verb|depot -> destination -> depot|. We
\textbf{centroid-snap} each sampled point to its containing BG and
look up the road km from $D_{\mathrm{road}}$. This is faster than
on-the-fly road queries per sample and accurate to within one
centroid-radius ($\approx 100$--$300$\,m for our 25-BG areas).

\paragraph{Summary of distance use per baseline.}
\begin{center}
\small
\begin{tabular}{l p{4.6cm} p{4.6cm}}
\toprule
Baseline & Optimisation distances & Evaluation distances \\
\midrule
Joint-Nom / Joint-DRO &
   BHH $\alpha_i = \beta\sqrt{\Lambda}\sum_b \sqrt{p_b A_b}$;
   $K_i$ via road km ($\verb|get_dist|$ override). &
   In-district NN-TSP on road km via
   $\verb|road_network.get_dist_matrix_km|$. \\
SP-Lit (Carlsson) &
   Same as Joint; ham-sandwich bisection on Euclidean grid in
   EPSG:2163 metres (BG centroids unchanged). &
   Same as Joint. \\
Multi-FR / Multi-FR-Detour &
   Angular partition + 2-opt + subpath LP on road-aware MDS
   positions $\Phi(b)$ (Euclidean over $\Phi$). &
   Custom Benders simulate; Euclidean over $\Phi$. \\
TP-Lit (Liu) &
   ADP+VRP cost matrices via Euclidean over $\Phi$
   ($n_{\mathrm{candidates}} = 0$ so candidates are real BG
   centroids). &
   Custom simulate; Euclidean over $\Phi$. \\
VCC &
   Dial-a-ride routing via Euclidean over $\Phi$. &
   Same. Continuous walk distances within blocks remain Euclidean. \\
OD &
   Single trip $\verb|depot -> dest -> depot|$. &
   Centroid-snap road km from $D_{\mathrm{road}}$. \\
\bottomrule
\end{tabular}
\end{center}

% -----------------------------------------------------------------------
\subsubsection{Demand}

Demand at each BG is proportional to its ACS 2023 total population
(table B01003).
Per-block probabilities are normalised to sum to one over the 25-BG
subset:

\begin{equation*}
  p_b = \frac{\mathrm{pop}_b}{\sum_{b' \in \mathrm{subset}} \mathrm{pop}_{b'}}.
\end{equation*}

Total subset populations vary across cities (Pittsburgh: 24{,}604;
Chicago: 63{,}522; Manhattan: 26{,}326; Los Angeles: 37{,}428),
which alone introduces meaningful per-block density differences ---
Manhattan's 25 BGs are roughly half the area of Pittsburgh's but
have similar total population. The system arrival rate $\Lambda$
is swept independently of population total; the population
distribution only sets the spatial pattern $p_b$.

% -----------------------------------------------------------------------
\subsubsection{Sweep grid (regime view)}

Identical to the synthetic $5\times5$ design (\S\ref{sec:regime-view})
so cross-instance comparison is direct:

\begin{itemize}
  \item $\Lambda \in \{50, 150, 300\}$ pax/h
  \item $w_r \in \{1, 3, 10\}$
  \item $w_v = 1.0$ (pinned)
  \item $c_f = 4.0$ \$/veh-h (5-yr fleet amortisation)
  \item $K = 3$ baseline (used by Multi-FR's fleet fall-back; partition methods sweep $K \in \{2,...,6\}$ internally)
\end{itemize}

For each $(\Lambda, w_r)$ point, every baseline runs its
internal hyperparameter sweep
(\S\ref{sec:sweep-params}: K for partition, $T_{\mathrm{fixed}}$ for
SP-Lit, fleet for OD/VCC, etc.) and reports the operating point
with the lowest aggregate cost
$C_{\mathrm{provider}} + w_r \cdot \bar{u} \cdot \Lambda$.

% -----------------------------------------------------------------------
\subsubsection{Reproducibility}

Three scripts implement the pipeline end to end:

\begin{enumerate}
  \item \texttt{scripts/experiments/build\_city\_geoid\_lists.py} \\
    Filter each state's TIGER BG shapefile by the city's focal bbox,
    keep the $25$ centermost BGs, save GEOID lists to
    \verb|data/city_bg_lists/<city>.txt|.
  \item \texttt{scripts/experiments/fetch\_city\_populations.py} \\
    Query the Census API for ACS 2023 B01003 per BG, save
    per-city JSON to \verb|data/city_bg_populations/<city>.json|.
  \item \texttt{scripts/experiments/run\_per\_city\_full.sh} \\
    Run the regime sweep across all 4 cities with road-aware
    optimisation, population-weighted demand, and all 7 baselines
    (Multi-FR + TP-Lit included). Resumable via per-setting JSON
    cache.
\end{enumerate}

The plotting helpers
\verb|scripts/plots/multi_city_designs_map.py| and the city-aware
extension of \verb|rebuild_latest_sweeps.py| produce the
4-panel scatter cloud and the $4 \times N_{\mathrm{baseline}}$
design map shown in Fig.~\ref{fig:cities-cloud} and
Fig.~\ref{fig:cities-designs}.

% -----------------------------------------------------------------------
\subsubsection{Known limitations}

\begin{itemize}
  \item \textbf{MDS approximation error.} Road-distance metrics are
    not exactly Euclidean-embeddable in 2D; pairwise distances
    after embedding can deviate by up to $\sim 30\%$ on irregular
    networks (Pittsburgh). A stricter alternative is to refactor
    each baseline's routing helpers to consume an explicit
    distance-matrix function; this is left for future work.
  \item \textbf{Continuous-point road distance in VCC.} VCC's
    inner dial-a-ride routing operates on continuous coordinates
    (interpolated meeting points). Those interior moves remain
    Euclidean over MDS positions; only depot, destination, and
    block-anchor distances are road-aware.
  \item \textbf{Demand temporality.} ACS B01003 is total residential
    population, not time-varying ridership. A more realistic
    deployment study would weight by trip-generation tables
    (e.g.\ ACS B08303 commuting) at peak hours. Population-weighted
    demand is a stand-in for cross-city methodology validation.
  \item \textbf{Single bbox per city.} Each city is represented by
    a single $\sim 5\,$km focal area. Multi-bbox runs (e.g.\
    Pittsburgh downtown + suburbs) would test whether topology
    rankings hold across intra-city heterogeneity.
\end{itemize}
 where
% needed. Self-contained: no other doc cross-references required.
% Requires: amsmath, booktabs, array, hyperref (optional).
% =============================================================================

\subsection{Real-City Experiment Design}\label{sec:real-city-design}

This section documents the design of the real-city baseline sweep
described in \S\ref{sec:multi-city}. The goal is to compare the seven
service-design baselines on the same regime grid as the synthetic
$5\times5$ run, but on \emph{actual} U.S.\ census block groups (BGs)
linked together by \emph{actual} street networks --- preserving
real demand density (population) and real driving distances.

% -----------------------------------------------------------------------
\subsubsection{Cities and their road topologies}

Four cities were chosen to span the major U.S.\ road-network archetypes:

\begin{center}
\small
\begin{tabular}{lll}
\toprule
City & Focal area / topology & BG suburb \\
\midrule
Pittsburgh, PA & Downtown + Oakland; hilly, river-cut, irregular grid &
   Allegheny Co.\ (FIPS 42003) \\
Chicago, IL    & Loop + River North; strict cardinal grid &
   Cook Co.\ (FIPS 17031) \\
Manhattan, NY  & Midtown; avenue-and-street grid plus Broadway &
   New York Co.\ (FIPS 36061) \\
Los Angeles, CA & DTLA; sprawled grid bisected by freeways &
   LA Co.\ (FIPS 06037) \\
\bottomrule
\end{tabular}
\end{center}

For each city we keep the $25$ block groups whose centroids are
closest to the centroid of a focal $\sim 5\,$km bbox. Twenty-five BGs
match the synthetic $5\times5$ instance and keep every per-baseline
sweep tractable on a workstation; the centermost selection avoids
boundary BGs whose road network is partially clipped.

% -----------------------------------------------------------------------
\subsubsection{Data sources}

\begin{description}
  \item[Block-group polygons] U.S.\ Census Bureau TIGER 2023
    shapefiles (one file per state, downloaded via
    \texttt{www2.census.gov/geo/tiger/TIGER2023/BG}). Reprojected to
    EPSG:2163 for area + Euclidean distance computations.
  \item[Population per BG] American Community Survey 5-year 2023
    table B01003 (total population), fetched per BG via the public
    Census Data API. Saved as a per-city JSON
    \verb|data/city_bg_populations/<city>.json|.
  \item[Road network] OpenStreetMap drive-network graph fetched via
    OSMnx 2.x for each city's bbox (with a $0.005^\circ$ pad).
    Cached on disk under \verb|cache/road_graph_*.pkl|.
\end{description}

The end-to-end fetch takes a few minutes per city; helper scripts
(\verb|build_city_geoid_lists.py|, \verb|fetch_city_populations.py|)
handle this once and write the artefacts to disk for reproducible
sweeps.

% -----------------------------------------------------------------------
\subsubsection{Distances: road network + MDS embedding}

The repository historically used Euclidean centroid-to-centroid
distances throughout. For real cities we replace these with
\textbf{road-network shortest-path distances} on three layers:

\paragraph{1. Centroid pair distances (\texttt{geodata.get\_dist}).}
After OSMnx fetches the drive graph, every BG centroid is mapped to
its nearest road-network node and an all-pairs shortest-path matrix
$D_{\mathrm{road}} \in \mathbb{R}^{N \times N}$ is precomputed
(\verb|GeoData.attach_road_distances|). Once attached,
\verb|GeoData.get_dist(b_1, b_2)| returns the road km from
$D_{\mathrm{road}}$, so any baseline that already calls
\verb|get_dist| (e.g.\ Joint-Nom/DRO and SP-Lit's linehaul
$K_i$) automatically becomes road-aware without code changes.

\paragraph{2. MDS-embedded positions for routing helpers.}
Several baselines (Multi-FR, TP-Lit, VCC) compute pairwise
distances via \verb|np.linalg.norm| over centroid coordinates
in EPSG:2163 metres. Replacing every call site with matrix
look-ups would be invasive. Instead, when the road-distance matrix
is attached we compute a \textbf{2D MDS embedding}
$\Phi : \mathrm{BG} \to \mathbb{R}^2$ such that
$\| \Phi(b_1) - \Phi(b_2) \|_2 \approx D_{\mathrm{road}}(b_1, b_2)$
and expose it as \verb|geodata.pos_road_km|. The embedding is
computed with scikit-learn's \verb|sklearn.manifold.MDS|
(precomputed dissimilarity, $n=2$, $4$ initialisations,
$300$ iterations).

A new helper \verb|_get_pos_road_aware(geodata, b)| returns
\verb|pos_road_km[b]| when available and falls back to the
Euclidean \verb|pos[b]| otherwise. We swapped this in for the
position lookups inside Multi-FR's line library, TP-Lit's ADP+VRP
candidate locations, and VCC's depot/dest/start positions. The
existing $L^2$-norm logic then operates on road-distance-consistent
geometry.

The MDS embedding is approximate --- general road-distance metrics
are not exactly Euclidean-embeddable in two dimensions --- but the
stress is moderate (typical pairwise error $\le 30\%$ on these 25-BG
instances) and the embedding propagates road awareness through
hundreds of $L^2$-norm call sites without bespoke refactors.

\paragraph{3. Continuous-sample distances in OD.}
OD samples customer destinations as continuous points within block
polygons, then routes \verb|depot -> destination -> depot|. We
\textbf{centroid-snap} each sampled point to its containing BG and
look up the road km from $D_{\mathrm{road}}$. This is faster than
on-the-fly road queries per sample and accurate to within one
centroid-radius ($\approx 100$--$300$\,m for our 25-BG areas).

\paragraph{Summary of distance use per baseline.}
\begin{center}
\small
\begin{tabular}{l p{4.6cm} p{4.6cm}}
\toprule
Baseline & Optimisation distances & Evaluation distances \\
\midrule
Joint-Nom / Joint-DRO &
   BHH $\alpha_i = \beta\sqrt{\Lambda}\sum_b \sqrt{p_b A_b}$;
   $K_i$ via road km ($\verb|get_dist|$ override). &
   In-district NN-TSP on road km via
   $\verb|road_network.get_dist_matrix_km|$. \\
SP-Lit (Carlsson) &
   Same as Joint; ham-sandwich bisection on Euclidean grid in
   EPSG:2163 metres (BG centroids unchanged). &
   Same as Joint. \\
Multi-FR / Multi-FR-Detour &
   Angular partition + 2-opt + subpath LP on road-aware MDS
   positions $\Phi(b)$ (Euclidean over $\Phi$). &
   Custom Benders simulate; Euclidean over $\Phi$. \\
TP-Lit (Liu) &
   ADP+VRP cost matrices via Euclidean over $\Phi$
   ($n_{\mathrm{candidates}} = 0$ so candidates are real BG
   centroids). &
   Custom simulate; Euclidean over $\Phi$. \\
VCC &
   Dial-a-ride routing via Euclidean over $\Phi$. &
   Same. Continuous walk distances within blocks remain Euclidean. \\
OD &
   Single trip $\verb|depot -> dest -> depot|$. &
   Centroid-snap road km from $D_{\mathrm{road}}$. \\
\bottomrule
\end{tabular}
\end{center}

% -----------------------------------------------------------------------
\subsubsection{Demand}

Demand at each BG is proportional to its ACS 2023 total population
(table B01003).
Per-block probabilities are normalised to sum to one over the 25-BG
subset:

\begin{equation*}
  p_b = \frac{\mathrm{pop}_b}{\sum_{b' \in \mathrm{subset}} \mathrm{pop}_{b'}}.
\end{equation*}

Total subset populations vary across cities (Pittsburgh: 24{,}604;
Chicago: 63{,}522; Manhattan: 26{,}326; Los Angeles: 37{,}428),
which alone introduces meaningful per-block density differences ---
Manhattan's 25 BGs are roughly half the area of Pittsburgh's but
have similar total population. The system arrival rate $\Lambda$
is swept independently of population total; the population
distribution only sets the spatial pattern $p_b$.

% -----------------------------------------------------------------------
\subsubsection{Sweep grid (regime view)}

Identical to the synthetic $5\times5$ design (\S\ref{sec:regime-view})
so cross-instance comparison is direct:

\begin{itemize}
  \item $\Lambda \in \{50, 150, 300\}$ pax/h
  \item $w_r \in \{1, 3, 10\}$
  \item $w_v = 1.0$ (pinned)
  \item $c_f = 4.0$ \$/veh-h (5-yr fleet amortisation)
  \item $K = 3$ baseline (used by Multi-FR's fleet fall-back; partition methods sweep $K \in \{2,...,6\}$ internally)
\end{itemize}

For each $(\Lambda, w_r)$ point, every baseline runs its
internal hyperparameter sweep
(\S\ref{sec:sweep-params}: K for partition, $T_{\mathrm{fixed}}$ for
SP-Lit, fleet for OD/VCC, etc.) and reports the operating point
with the lowest aggregate cost
$C_{\mathrm{provider}} + w_r \cdot \bar{u} \cdot \Lambda$.

% -----------------------------------------------------------------------
\subsubsection{Reproducibility}

Three scripts implement the pipeline end to end:

\begin{enumerate}
  \item \texttt{scripts/experiments/build\_city\_geoid\_lists.py} \\
    Filter each state's TIGER BG shapefile by the city's focal bbox,
    keep the $25$ centermost BGs, save GEOID lists to
    \verb|data/city_bg_lists/<city>.txt|.
  \item \texttt{scripts/experiments/fetch\_city\_populations.py} \\
    Query the Census API for ACS 2023 B01003 per BG, save
    per-city JSON to \verb|data/city_bg_populations/<city>.json|.
  \item \texttt{scripts/experiments/run\_per\_city\_full.sh} \\
    Run the regime sweep across all 4 cities with road-aware
    optimisation, population-weighted demand, and all 7 baselines
    (Multi-FR + TP-Lit included). Resumable via per-setting JSON
    cache.
\end{enumerate}

The plotting helpers
\verb|scripts/plots/multi_city_designs_map.py| and the city-aware
extension of \verb|rebuild_latest_sweeps.py| produce the
4-panel scatter cloud and the $4 \times N_{\mathrm{baseline}}$
design map shown in Fig.~\ref{fig:cities-cloud} and
Fig.~\ref{fig:cities-designs}.

% -----------------------------------------------------------------------
\subsubsection{Known limitations}

\begin{itemize}
  \item \textbf{MDS approximation error.} Road-distance metrics are
    not exactly Euclidean-embeddable in 2D; pairwise distances
    after embedding can deviate by up to $\sim 30\%$ on irregular
    networks (Pittsburgh). A stricter alternative is to refactor
    each baseline's routing helpers to consume an explicit
    distance-matrix function; this is left for future work.
  \item \textbf{Continuous-point road distance in VCC.} VCC's
    inner dial-a-ride routing operates on continuous coordinates
    (interpolated meeting points). Those interior moves remain
    Euclidean over MDS positions; only depot, destination, and
    block-anchor distances are road-aware.
  \item \textbf{Demand temporality.} ACS B01003 is total residential
    population, not time-varying ridership. A more realistic
    deployment study would weight by trip-generation tables
    (e.g.\ ACS B08303 commuting) at peak hours. Population-weighted
    demand is a stand-in for cross-city methodology validation.
  \item \textbf{Single bbox per city.} Each city is represented by
    a single $\sim 5\,$km focal area. Multi-bbox runs (e.g.\
    Pittsburgh downtown + suburbs) would test whether topology
    rankings hold across intra-city heterogeneity.
\end{itemize}
